%!TEX program = xelatex
\documentclass[12pt,a4paper,oneside]{article}

% ============================================================================
% PACKAGES & CONFIGURATION
% ============================================================================
\usepackage[utf-8]{inputenc}
\usepackage[left=2.5cm,right=2.5cm,top=3cm,bottom=3cm]{geometry}
\usepackage{graphicx}
\usepackage{fancyhdr}
\usepackage{color}
\usepackage{listings}
\usepackage{hyperref}
\usepackage{xcolor}
\usepackage{tcolorbox}
\usepackage{array}
\usepackage{booktabs}
\usepackage{amsmath}
\usepackage{amssymb}
\usepackage{float}
\usepackage{titlesec}
\usepackage{enumitem}
\usepackage{fancyvrb}

% Color definitions
\definecolor{codegreen}{rgb}{0.0, 0.5, 0.0}
\definecolor{codegray}{rgb}{0.5, 0.5, 0.5}
\definecolor{codepurple}{rgb}{0.58, 0, 0.82}
\definecolor{backcolour}{rgb}{0.95, 0.95, 0.95}
\definecolor{lightblue}{rgb}{0.9, 0.95, 1.0}
\definecolor{darkblue}{rgb}{0.0, 0.2, 0.6}
\definecolor{lightyellow}{rgb}{1.0, 0.98, 0.8}
\definecolor{orange}{rgb}{1.0, 0.65, 0.0}

% Code listing configuration
\lstset{
    backgroundcolor=\color{backcolour},
    commentstyle=\color{codegreen},
    keywordstyle=\color{codepurple},
    numberstyle=\tiny\color{codegray},
    stringstyle=\color{codegreen},
    basicstyle=\ttfamily\small,
    breakatwhitespace=false,
    breaklines=true,
    captionpos=b,
    keepspaces=true,
    numbers=left,
    numbersep=5pt,
    showspaces=false,
    showstringspaces=false,
    showtabs=false,
    tabsize=2,
    frame=single,
    language=javascript,
}

% Header and footer
\pagestyle{fancy}
\fancyhf{}
\rhead{\textcolor{darkblue}{\textit{VisioLearn Frontend Guide}}}
\lhead{\textcolor{darkblue}{Frontend Dev}}
\cfoot{\thepage}
\renewcommand{\headrulewidth}{0.5pt}
\renewcommand{\footrulewidth}{0.5pt}

% Section styling
\titleformat{\section}{\Large\bfseries\color{darkblue}}{}{0pt}{}[\titlerule]
\titleformat{\subsection}{\large\bfseries\color{orange}}{}{0pt}{}
\titleformat{\subsubsection}{\bfseries\color{darkblue}}{}{0pt}{}

% ============================================================================
% TITLE PAGE
% ============================================================================
\title{%
    \vspace{2cm}
    \Huge\textbf{\textcolor{darkblue}{VisioLearn Frontend}} \\
    \vspace{0.5cm}
    \Large\textit{UI/UX Development Guide} \\
    \vspace{2cm}
    \large\textbf{Complete Frontend Implementation Reference}
}

\author{%
    Frontend Development Team \\
    \vspace{0.5cm}
    Integrated with Backend API at \\
    \href{https://visiolearn-backend.onrender.com}{https://visiolearn-backend.onrender.com}
}

\date{\vspace{1cm}\large \textbf{Version 1.0} \\ May 2026}

% ============================================================================
% DOCUMENT
% ============================================================================
\begin{document}

% Title page
\maketitle
\thispagestyle{empty}
\newpage

% ============================================================================
% TABLE OF CONTENTS
% ============================================================================
\tableofcontents
\newpage

% ============================================================================
% SECTION 1: OVERVIEW
% ============================================================================
\section{Frontend Overview}

VisioLearn is a web-based UI for an AI-powered audio learning platform. The frontend enables users to:

\begin{itemize}
    \item \textbf{Authenticate}: Login, manage sessions, logout securely
    \item \textbf{Upload Content}: Teachers upload lesson materials (PDF, DOCX, TXT)
    \item \textbf{Voice Learning}: Students interact via voice commands
    \item \textbf{Track Progress}: View analytics and learning history
    \item \textbf{Manage Users}: Administrative user management
\end{itemize}

\subsection{Tech Stack Recommendations}

\begin{table}[H]
\centering
\small
\begin{tabular}{|l|l|l|}
\hline
\textbf{Layer} & \textbf{Recommended} & \textbf{Alternative} \\
\hline
Framework & React 18+ & Vue 3, Angular \\
State Mgmt & Redux Toolkit & Zustand, Jotai \\
HTTP Client & Axios & Fetch API, React Query \\
UI Library & Material-UI & Ant Design, Chakra \\
Build Tool & Vite & Webpack, Next.js \\
Styling & Tailwind CSS & Emotion, Styled Components \\
Testing & Jest + RTL & Vitest, Playwright \\
\hline
\end{tabular}
\caption{Recommended Frontend Technology Stack}
\end{table}

\subsection{Backend API Base URL}

\begin{tcolorbox}[colback=lightblue, colframe=darkblue]
\texttt{https://visiolearn-backend.onrender.com/api/v1}
\end{tcolorbox}

% ============================================================================
% SECTION 2: AUTHENTICATION FLOW
% ============================================================================
\section{Authentication \& Authorization}

\subsection{Login Flow (UX)}

The authentication flow involves three main steps:

\begin{enumerate}
    \item \textbf{User Input}: Email and password form submission
    \item \textbf{API Request}: POST to \texttt{/auth/login}
    \item \textbf{Token Storage}: Store access and refresh tokens
    \item \textbf{Redirect}: Route to appropriate dashboard
\end{enumerate}

\subsection{Login Component Implementation}

\begin{lstlisting}[language=javascript]
// pages/Login.jsx
import { useState } from 'react';
import { useNavigate } from 'react-router-dom';
import { useDispatch } from 'react-redux';
import { setAuthTokens, setUser } from '../store/authSlice';
import api from '../api/client';

export default function Login() {
  const [email, setEmail] = useState('');
  const [password, setPassword] = useState('');
  const [loading, setLoading] = useState(false);
  const [error, setError] = useState('');
  const navigate = useNavigate();
  const dispatch = useDispatch();

  const handleLogin = async (e) => {
    e.preventDefault();
    setLoading(true);
    setError('');

    try {
      const response = await api.post('/auth/login', {
        email: email.toLowerCase(),
        password
      });

      // Store tokens
      dispatch(setAuthTokens({
        accessToken: response.data.access_token,
        refreshToken: response.data.refresh_token
      }));

      // Redirect based on user role
      const decodedToken = jwtDecode(response.data.access_token);
      const userRole = decodedToken.role;

      if (userRole === 'admin') navigate('/admin-dashboard');
      else if (userRole === 'teacher') navigate('/teacher-dashboard');
      else if (userRole === 'student') navigate('/student-dashboard');

    } catch (err) {
      setError(
        err.response?.data?.detail || 
        'Login failed. Please try again.'
      );
    } finally {
      setLoading(false);
    }
  };

  return (
    <div className="login-container">
      <form onSubmit={handleLogin}>
        <h1>VisioLearn Login</h1>
        
        <input
          type="email"
          value={email}
          onChange={(e) => setEmail(e.target.value)}
          placeholder="Email"
          required
        />
        
        <input
          type="password"
          value={password}
          onChange={(e) => setPassword(e.target.value)}
          placeholder="Password"
          required
        />
        
        {error && <div className="error">{error}</div>}
        
        <button type="submit" disabled={loading}>
          {loading ? 'Logging in...' : 'Login'}
        </button>
      </form>
    </div>
  );
}
\end{lstlisting}

\subsection{Token Management}

\textbf{Access Token Storage Strategy:}

\begin{itemize}
    \item \textbf{Access Token}: Store in Redux state (NOT localStorage for XSS protection)
    \item \textbf{Refresh Token}: Store in httpOnly cookie (handled by backend)
    \item \textbf{Expiry}: 60 minutes - implement refresh before expiry
    \item \textbf{Logout}: Clear tokens on logout
\end{itemize}

\subsection{Token Refresh Implementation}

\begin{lstlisting}[language=javascript]
// api/interceptor.js
import api from './client';
import store from '../store';
import { setAuthTokens, clearAuth } from '../store/authSlice';

// Add response interceptor for token refresh
api.interceptors.response.use(
  response => response,
  async error => {
    const originalRequest = error.config;

    // If 401 and not already retried
    if (error.response?.status === 401 && 
        !originalRequest._retry) {
      
      originalRequest._retry = true;
      const state = store.getState();
      const refreshToken = state.auth.refreshToken;

      if (!refreshToken) {
        store.dispatch(clearAuth());
        window.location.href = '/login';
        return Promise.reject(error);
      }

      try {
        // Request new tokens
        const response = await api.post('/auth/refresh', {
          refresh_token: refreshToken
        });

        // Update tokens in store
        store.dispatch(setAuthTokens({
          accessToken: response.data.access_token,
          refreshToken: response.data.refresh_token
        }));

        // Retry original request with new token
        originalRequest.headers.Authorization = 
          `Bearer ${response.data.access_token}`;
        return api(originalRequest);

      } catch (refreshError) {
        store.dispatch(clearAuth());
        window.location.href = '/login';
        return Promise.reject(refreshError);
      }
    }

    return Promise.reject(error);
  }
);
\end{lstlisting}

% ============================================================================
% SECTION 3: API CLIENT SETUP
% ============================================================================
\section{API Client Configuration}

\subsection{Axios Configuration}

\begin{lstlisting}[language=javascript]
// api/client.js
import axios from 'axios';
import store from '../store';

const api = axios.create({
  baseURL: process.env.REACT_APP_API_BASE_URL || 
           'https://visiolearn-backend.onrender.com/api/v1',
  timeout: 30000,
  headers: {
    'Content-Type': 'application/json',
    'Accept': 'application/json'
  }
});

// Request interceptor to add auth token
api.interceptors.request.use(
  config => {
    const state = store.getState();
    const token = state.auth.accessToken;
    
    if (token) {
      config.headers.Authorization = `Bearer ${token}`;
    }
    
    return config;
  },
  error => Promise.reject(error)
);

export default api;
\end{lstlisting}

\subsection{Environment Variables}

Create \texttt{.env} file:

\begin{lstlisting}
REACT_APP_API_BASE_URL=https://visiolearn-backend.onrender.com/api/v1
REACT_APP_API_TIMEOUT=30000
REACT_APP_ENV=production
\end{lstlisting}

% ============================================================================
% SECTION 4: COMPONENT ARCHITECTURE
% ============================================================================
\section{Component Architecture}

\subsection{Directory Structure}

\begin{lstlisting}
frontend/
├── src/
│   ├── components/
│   │   ├── Auth/
│   │   │   ├── Login.jsx
│   │   │   ├── Logout.jsx
│   │   │   └── ProtectedRoute.jsx
│   │   ├── LessonNotes/
│   │   │   ├── NoteUpload.jsx
│   │   │   ├── NoteList.jsx
│   │   │   └── NoteDetail.jsx
│   │   ├── VoiceSessions/
│   │   │   ├── SessionStart.jsx
│   │   │   ├── SessionActive.jsx
│   │   │   ├── SessionEnd.jsx
│   │   │   └── InteractionLog.jsx
│   │   └── Common/
│   │       ├── Header.jsx
│   │       ├── Sidebar.jsx
│   │       └── LoadingSpinner.jsx
│   ├── pages/
│   │   ├── AdminDashboard.jsx
│   │   ├── TeacherDashboard.jsx
│   │   ├── StudentDashboard.jsx
│   │   └── NotFound.jsx
│   ├── store/
│   │   ├── authSlice.js
│   │   ├── notesSlice.js
│   │   ├── voiceSlice.js
│   │   └── index.js
│   ├── hooks/
│   │   ├── useAuth.js
│   │   ├── useNotes.js
│   │   └── useVoiceSession.js
│   ├── api/
│   │   ├── client.js
│   │   ├── interceptor.js
│   │   └── endpoints.js
│   ├── utils/
│   │   ├── formatting.js
│   │   ├── validation.js
│   │   └── errorHandler.js
│   ├── styles/
│   │   ├── variables.css
│   │   ├── globals.css
│   │   └── components.css
│   ├── App.jsx
│   └── index.jsx
└── package.json
\end{lstlisting}

\subsection{Protected Route Component}

\begin{lstlisting}[language=javascript]
// components/Auth/ProtectedRoute.jsx
import { Navigate } from 'react-router-dom';
import { useSelector } from 'react-redux';

export default function ProtectedRoute({ 
  children, 
  requiredRole 
}) {
  const { accessToken, user } = useSelector(
    state => state.auth
  );

  if (!accessToken) {
    return <Navigate to="/login" replace />;
  }

  if (requiredRole && user?.role !== requiredRole) {
    return <Navigate to="/unauthorized" replace />;
  }

  return children;
}

// Usage:
// <ProtectedRoute requiredRole="teacher">
//   <TeacherDashboard />
// </ProtectedRoute>
\end{lstlisting}

% ============================================================================
% SECTION 5: LESSON NOTES UI
% ============================================================================
\section{Lesson Notes Feature}

\subsection{Note Upload Component}

\begin{lstlisting}[language=javascript]
// components/LessonNotes/NoteUpload.jsx
import { useState } from 'react';
import { useDispatch } from 'react-redux';
import { uploadNote } from '../store/notesSlice';
import api from '../api/client';

const SUPPORTED_TYPES = [
  'application/pdf',
  'application/vnd.openxmlformats-officedocument.wordprocessingml.document',
  'text/plain'
];

export default function NoteUpload() {
  const [formData, setFormData] = useState({
    title: '',
    subject: '',
    gradeLevel: '',
    file: null
  });
  const [loading, setLoading] = useState(false);
  const [error, setError] = useState('');
  const [success, setSuccess] = useState('');
  const dispatch = useDispatch();

  const handleFileChange = (e) => {
    const file = e.target.files[0];
    
    // Validate file type
    if (!SUPPORTED_TYPES.includes(file.type)) {
      setError('Unsupported file type. Use PDF, DOCX, or TXT');
      return;
    }

    // Validate file size (50MB max)
    if (file.size > 50 * 1024 * 1024) {
      setError('File too large. Maximum 50MB allowed');
      return;
    }

    setFormData(prev => ({ ...prev, file }));
    setError('');
  };

  const handleSubmit = async (e) => {
    e.preventDefault();
    setLoading(true);
    setError('');

    try {
      const formDataToSend = new FormData();
      formDataToSend.append('file', formData.file);
      formDataToSend.append('title', formData.title);
      formDataToSend.append('subject', formData.subject);
      formDataToSend.append('grade_level', formData.gradeLevel);

      const response = await api.post(
        '/notes/upload',
        formDataToSend,
        {
          headers: {
            'Content-Type': 'multipart/form-data'
          }
        }
      );

      setSuccess(
        `Successfully uploaded: ${formData.title}`
      );
      dispatch(uploadNote(response.data));
      
      // Reset form
      setFormData({
        title: '',
        subject: '',
        gradeLevel: '',
        file: null
      });

    } catch (err) {
      setError(
        err.response?.data?.detail || 
        'Upload failed. Try again.'
      );
    } finally {
      setLoading(false);
    }
  };

  return (
    <div className="upload-container">
      <h2>Upload Lesson Note</h2>
      
      <form onSubmit={handleSubmit}>
        <div className="form-group">
          <label>Title</label>
          <input
            type="text"
            value={formData.title}
            onChange={(e) => setFormData(
              prev => ({ ...prev, title: e.target.value })
            )}
            required
          />
        </div>

        <div className="form-group">
          <label>Subject *</label>
          <input
            type="text"
            value={formData.subject}
            onChange={(e) => setFormData(
              prev => ({ ...prev, subject: e.target.value })
            )}
            required
          />
        </div>

        <div className="form-group">
          <label>Grade Level *</label>
          <select
            value={formData.gradeLevel}
            onChange={(e) => setFormData(
              prev => ({ ...prev, gradeLevel: e.target.value })
            )}
            required
          >
            <option value="">Select grade level</option>
            <option value="8">Grade 8</option>
            <option value="9">Grade 9</option>
            <option value="10">Grade 10</option>
            <option value="11">Grade 11</option>
            <option value="12">Grade 12</option>
          </select>
        </div>

        <div className="form-group">
          <label>File (PDF, DOCX, TXT) *</label>
          <input
            type="file"
            onChange={handleFileChange}
            accept=".pdf,.docx,.txt"
            required
          />
          {formData.file && (
            <p className="file-info">
              Selected: {formData.file.name}
            </p>
          )}
        </div>

        {error && <div className="alert alert-error">{error}</div>}
        {success && <div className="alert alert-success">{success}</div>}

        <button type="submit" disabled={loading}>
          {loading ? 'Uploading...' : 'Upload Note'}
        </button>
      </form>
    </div>
  );
}
\end{lstlisting}

\subsection{Progress Indicators}

Display upload progress for files:

\begin{lstlisting}[language=javascript]
// Show progress during upload
const config = {
  onUploadProgress: (progressEvent) => {
    const percentComplete = Math.round(
      (progressEvent.loaded * 100) / progressEvent.total
    );
    setUploadProgress(percentComplete);
  }
};

const response = await api.post(
  '/notes/upload',
  formDataToSend,
  config
);
\end{lstlisting}

% ============================================================================
% SECTION 6: VOICE SESSIONS UI
% ============================================================================
\section{Voice Sessions Feature}

\subsection{Voice Session Flow}

The voice session has three states:

\begin{enumerate}
    \item \textbf{Start Session}: Initialize before learning
    \item \textbf{Active Session}: Log interactions during learning
    \item \textbf{End Session}: Complete and get summary
\end{enumerate}

\subsection{Start Session Component}

\begin{lstlisting}[language=javascript]
// components/VoiceSessions/SessionStart.jsx
import { useState, useEffect } from 'react';
import { useNavigate } from 'react-router-dom';
import { useDispatch } from 'react-redux';
import { startSession } from '../store/voiceSlice';
import api from '../api/client';

export default function SessionStart() {
  const [students, setStudents] = useState([]);
  const [notes, setNotes] = useState([]);
  const [units, setUnits] = useState([]);
  const [selectedStudent, setSelectedStudent] = useState('');
  const [selectedNote, setSelectedNote] = useState('');
  const [selectedUnit, setSelectedUnit] = useState('');
  const [loading, setLoading] = useState(false);
  const navigate = useNavigate();
  const dispatch = useDispatch();

  useEffect(() => {
    // Fetch available students (for admin/teachers)
    fetchStudents();
    // Fetch available notes
    fetchNotes();
  }, []);

  useEffect(() => {
    // Fetch units when note is selected
    if (selectedNote) {
      fetchUnits(selectedNote);
    }
  }, [selectedNote]);

  const fetchStudents = async () => {
    try {
      const response = await api.get('/users/students');
      setStudents(response.data);
    } catch (error) {
      console.error('Failed to fetch students:', error);
    }
  };

  const fetchNotes = async () => {
    try {
      const response = await api.get('/notes');
      setNotes(response.data);
    } catch (error) {
      console.error('Failed to fetch notes:', error);
    }
  };

  const fetchUnits = async (noteId) => {
    try {
      const response = await api.get(`/notes/${noteId}/units`);
      setUnits(response.data);
    } catch (error) {
      console.error('Failed to fetch units:', error);
    }
  };

  const handleStartSession = async () => {
    setLoading(true);

    try {
      const response = await api.post('/voice/session/start', {
        student_id: selectedStudent,
        note_id: selectedNote,
        unit_id: selectedUnit
      });

      dispatch(startSession({
        sessionId: response.data.session_id,
        studentId: response.data.student_id,
        noteId: response.data.note_id,
        unitId: response.data.unit_id,
        startTime: response.data.started_at
      }));

      navigate(`/voice/session/${response.data.session_id}`);

    } catch (error) {
      alert('Failed to start session: ' + 
            error.response?.data?.detail);
    } finally {
      setLoading(false);
    }
  };

  return (
    <div className="session-start-container">
      <h1>Start Voice Session</h1>
      
      <div className="form-group">
        <label>Select Student</label>
        <select
          value={selectedStudent}
          onChange={(e) => setSelectedStudent(e.target.value)}
        >
          <option value="">Choose a student</option>
          {students.map(student => (
            <option key={student.id} value={student.id}>
              {student.full_name}
            </option>
          ))}
        </select>
      </div>

      <div className="form-group">
        <label>Select Lesson</label>
        <select
          value={selectedNote}
          onChange={(e) => setSelectedNote(e.target.value)}
        >
          <option value="">Choose a lesson</option>
          {notes.map(note => (
            <option key={note.id} value={note.id}>
              {note.title}
            </option>
          ))}
        </select>
      </div>

      <div className="form-group">
        <label>Select Unit</label>
        <select
          value={selectedUnit}
          onChange={(e) => setSelectedUnit(e.target.value)}
          disabled={!selectedNote}
        >
          <option value="">Choose a unit</option>
          {units.map(unit => (
            <option key={unit.id} value={unit.id}>
              {unit.title}
            </option>
          ))}
        </select>
      </div>

      <button
        onClick={handleStartSession}
        disabled={!selectedStudent || !selectedNote || 
                  !selectedUnit || loading}
        className="btn-primary"
      >
        {loading ? 'Starting...' : 'Start Session'}
      </button>
    </div>
  );
}
\end{lstlisting}

\subsection{Active Session Component}

\begin{lstlisting}[language=javascript]
// components/VoiceSessions/SessionActive.jsx
import { useState, useEffect } from 'react';
import { useParams } from 'react-router-dom';
import { useDispatch, useSelector } from 'react-redux';
import { logInteraction } from '../store/voiceSlice';
import api from '../api/client';

export default function SessionActive() {
  const { sessionId } = useParams();
  const [isListening, setIsListening] = useState(false);
  const [interactions, setInteractions] = useState([]);
  const [sessionTime, setSessionTime] = useState(0);
  const dispatch = useDispatch();

  useEffect(() => {
    // Start timer
    const timer = setInterval(() => {
      setSessionTime(t => t + 1);
    }, 1000);

    // Load existing interactions
    loadInteractions();

    return () => clearInterval(timer);
  }, [sessionId]);

  const loadInteractions = async () => {
    try {
      const response = await api.get(
        `/voice/session/${sessionId}/interactions`
      );
      setInteractions(response.data);
    } catch (error) {
      console.error('Failed to load interactions:', error);
    }
  };

  const recordInteraction = async (type, command) => {
    try {
      const response = await api.post(
        '/voice/session/event',
        {
          session_id: sessionId,
          interaction_type: type,
          command: command,
          confidence: 0.95,
          response: 'Processing...'
        }
      );

      setInteractions([...interactions, response.data]);
      dispatch(logInteraction(response.data));

    } catch (error) {
      alert('Failed to log interaction: ' + 
            error.response?.data?.detail);
    }
  };

  const formatTime = (seconds) => {
    const mins = Math.floor(seconds / 60);
    const secs = seconds % 60;
    return `${mins}:${secs.toString().padStart(2, '0')}`;
  };

  return (
    <div className="session-active-container">
      <div className="session-header">
        <h1>Voice Session Active</h1>
        <div className="timer">
          Time: {formatTime(sessionTime)}
        </div>
      </div>

      <div className="interaction-controls">
        <button 
          className={`btn ${isListening ? 'recording' : ''}`}
          onClick={() => setIsListening(!isListening)}
        >
          {isListening ? '🎤 Recording...' : '🎤 Start Recording'}
        </button>

        <button 
          onClick={() => recordInteraction('answer', 'Sample answer')}
          className="btn-secondary"
        >
          Log Answer
        </button>

        <button 
          onClick={() => recordInteraction('next', 'Next unit')}
          className="btn-secondary"
        >
          Next Unit
        </button>

        <button 
          onClick={() => recordInteraction('pause', 'Pause')}
          className="btn-warning"
        >
          Pause
        </button>
      </div>

      <div className="interactions-log">
        <h3>Interaction Log</h3>
        {interactions.length === 0 ? (
          <p>No interactions yet</p>
        ) : (
          <ul>
            {interactions.map((interaction, idx) => (
              <li key={idx}>
                <strong>{interaction.detected_intent}:</strong> 
                {interaction.student_transcript}
              </li>
            ))}
          </ul>
        )}
      </div>
    </div>
  );
}
\end{lstlisting}

\subsection{End Session Component}

\begin{lstlisting}[language=javascript]
// components/VoiceSessions/SessionEnd.jsx
import { useState } from 'react';
import { useParams, useNavigate } from 'react-router-dom';
import api from '../api/client';

export default function SessionEnd() {
  const { sessionId } = useParams();
  const [formData, setFormData] = useState({
    durationSeconds: 0,
    questionsAnswered: 0,
    totalScore: 0
  });
  const [result, setResult] = useState(null);
  const navigate = useNavigate();

  const handleEndSession = async () => {
    try {
      const response = await api.post(
        '/voice/session/end',
        {
          session_id: sessionId,
          duration_seconds: formData.durationSeconds,
          questions_answered: formData.questionsAnswered,
          total_score: formData.totalScore
        }
      );

      setResult(response.data);

    } catch (error) {
      alert('Failed to end session: ' + 
            error.response?.data?.detail);
    }
  };

  if (result) {
    return (
      <div className="session-result">
        <h1>Session Complete!</h1>
        <div className="results-card">
          <p><strong>Duration:</strong> 
             {Math.floor(result.duration_seconds / 60)} minutes
          </p>
          <p><strong>Questions Answered:</strong> 
             {result.questions_answered}
          </p>
          <p><strong>Average Score:</strong> 
             {(result.average_score * 100).toFixed(1)}%
          </p>
        </div>
        <button 
          onClick={() => navigate('/dashboard')}
          className="btn-primary"
        >
          Back to Dashboard
        </button>
      </div>
    );
  }

  return (
    <div className="session-end-container">
      <h1>End Voice Session</h1>
      
      <div className="form-group">
        <label>Duration (seconds)</label>
        <input
          type="number"
          value={formData.durationSeconds}
          onChange={(e) => setFormData(prev => ({
            ...prev,
            durationSeconds: parseInt(e.target.value)
          }))}
        />
      </div>

      <div className="form-group">
        <label>Questions Answered</label>
        <input
          type="number"
          value={formData.questionsAnswered}
          onChange={(e) => setFormData(prev => ({
            ...prev,
            questionsAnswered: parseInt(e.target.value)
          }))}
        />
      </div>

      <div className="form-group">
        <label>Total Score</label>
        <input
          type="number"
          step="0.1"
          value={formData.totalScore}
          onChange={(e) => setFormData(prev => ({
            ...prev,
            totalScore: parseFloat(e.target.value)
          }))}
        />
      </div>

      <button 
        onClick={handleEndSession}
        className="btn-primary"
      >
        End Session
      </button>
    </div>
  );
}
\end{lstlisting}

% ============================================================================
% SECTION 7: STATE MANAGEMENT
% ============================================================================
\section{State Management with Redux}

\subsection{Auth Slice}

\begin{lstlisting}[language=javascript]
// store/authSlice.js
import { createSlice } from '@reduxjs/toolkit';
import jwtDecode from 'jwt-decode';

const initialState = {
  accessToken: null,
  refreshToken: null,
  user: null,
  isAuthenticated: false,
  loading: false,
  error: null
};

const authSlice = createSlice({
  name: 'auth',
  initialState,
  reducers: {
    setAuthTokens: (state, action) => {
      state.accessToken = action.payload.accessToken;
      state.refreshToken = action.payload.refreshToken;
      
      // Decode token to get user info
      const decoded = jwtDecode(action.payload.accessToken);
      state.user = {
        id: decoded.sub,
        role: decoded.role
      };
      state.isAuthenticated = true;
    },
    clearAuth: (state) => {
      state.accessToken = null;
      state.refreshToken = null;
      state.user = null;
      state.isAuthenticated = false;
    },
    setLoading: (state, action) => {
      state.loading = action.payload;
    },
    setError: (state, action) => {
      state.error = action.payload;
    }
  }
});

export const { 
  setAuthTokens, 
  clearAuth, 
  setLoading, 
  setError 
} = authSlice.actions;

export default authSlice.reducer;
\end{lstlisting}

\subsection{Voice Session Slice}

\begin{lstlisting}[language=javascript]
// store/voiceSlice.js
import { createSlice } from '@reduxjs/toolkit';

const initialState = {
  currentSession: null,
  sessions: [],
  interactions: [],
  loading: false,
  error: null
};

const voiceSlice = createSlice({
  name: 'voice',
  initialState,
  reducers: {
    startSession: (state, action) => {
      state.currentSession = action.payload;
      state.interactions = [];
    },
    logInteraction: (state, action) => {
      state.interactions.push(action.payload);
    },
    endSession: (state, action) => {
      state.currentSession = null;
      state.sessions.push(action.payload);
    },
    setLoading: (state, action) => {
      state.loading = action.payload;
    },
    setError: (state, action) => {
      state.error = action.payload;
    }
  }
});

export const { 
  startSession, 
  logInteraction, 
  endSession 
} = voiceSlice.actions;

export default voiceSlice.reducer;
\end{lstlisting}

% ============================================================================
% SECTION 8: CUSTOM HOOKS
% ============================================================================
\section{Custom Hooks}

\subsection{useAuth Hook}

\begin{lstlisting}[language=javascript]
// hooks/useAuth.js
import { useSelector, useDispatch } from 'react-redux';
import { 
  setAuthTokens, 
  clearAuth 
} from '../store/authSlice';

export function useAuth() {
  const dispatch = useDispatch();
  const auth = useSelector(state => state.auth);

  return {
    ...auth,
    logout: () => dispatch(clearAuth()),
    setTokens: (tokens) => 
      dispatch(setAuthTokens(tokens))
  };
}
\end{lstlisting}

\subsection{useVoiceSession Hook}

\begin{lstlisting}[language=javascript]
// hooks/useVoiceSession.js
import { useSelector, useDispatch } from 'react-redux';
import { 
  startSession, 
  logInteraction, 
  endSession 
} from '../store/voiceSlice';
import api from '../api/client';

export function useVoiceSession() {
  const dispatch = useDispatch();
  const voice = useSelector(state => state.voice);

  const start = async (studentId, noteId, unitId) => {
    try {
      const response = await api.post(
        '/voice/session/start',
        { student_id: studentId, note_id: noteId, 
          unit_id: unitId }
      );
      dispatch(startSession(response.data));
      return response.data;
    } catch (error) {
      throw new Error(
        error.response?.data?.detail || 
        'Failed to start session'
      );
    }
  };

  const log = async (sessionId, type, command) => {
    try {
      const response = await api.post(
        '/voice/session/event',
        {
          session_id: sessionId,
          interaction_type: type,
          command: command
        }
      );
      dispatch(logInteraction(response.data));
      return response.data;
    } catch (error) {
      throw new Error(
        error.response?.data?.detail || 
        'Failed to log interaction'
      );
    }
  };

  const end = async (sessionId, duration, 
                      questionsAnswered, totalScore) => {
    try {
      const response = await api.post(
        '/voice/session/end',
        {
          session_id: sessionId,
          duration_seconds: duration,
          questions_answered: questionsAnswered,
          total_score: totalScore
        }
      );
      dispatch(endSession(response.data));
      return response.data;
    } catch (error) {
      throw new Error(
        error.response?.data?.detail || 
        'Failed to end session'
      );
    }
  };

  return {
    ...voice,
    start,
    log,
    end
  };
}
\end{lstlisting}

% ============================================================================
% SECTION 9: ERROR HANDLING & VALIDATION
% ============================================================================
\section{Error Handling \& Validation}

\subsection{Error Handler Utility}

\begin{lstlisting}[language=javascript]
// utils/errorHandler.js
export const handleApiError = (error) => {
  if (error.response) {
    // API responded with error status
    const status = error.response.status;
    const detail = error.response.data?.detail;

    switch (status) {
      case 401:
        return 'Please log in again';
      case 403:
        return 'You do not have permission';
      case 404:
        return 'Resource not found';
      case 400:
        return detail || 'Invalid request';
      case 500:
        return 'Server error. Try again later';
      default:
        return detail || 'An error occurred';
    }
  } else if (error.request) {
    return 'No response from server';
  }
  return error.message;
};

export const validateEmail = (email) => {
  const regex = /^[^\s@]+@[^\s@]+\.[^\s@]+$/;
  return regex.test(email);
};

export const validatePassword = (password) => {
  return password.length >= 8;
};
\end{lstlisting}

\subsection{Form Validation Hook}

\begin{lstlisting}[language=javascript]
// hooks/useForm.js
import { useState } from 'react';

export function useForm(initialValues, onSubmit) {
  const [values, setValues] = useState(initialValues);
  const [errors, setErrors] = useState({});
  const [touched, setTouched] = useState({});

  const handleChange = (e) => {
    const { name, value } = e.target;
    setValues(prev => ({ ...prev, [name]: value }));
  };

  const handleBlur = (e) => {
    const { name } = e.target;
    setTouched(prev => ({ ...prev, [name]: true }));
  };

  const handleSubmit = async (e) => {
    e.preventDefault();
    await onSubmit(values);
  };

  return {
    values,
    errors,
    touched,
    handleChange,
    handleBlur,
    handleSubmit,
    setErrors,
    setValues
  };
}
\end{lstlisting}

% ============================================================================
% SECTION 10: TESTING
% ============================================================================
\section{Frontend Testing}

\subsection{Test Structure}

\begin{lstlisting}[language=javascript]
// components/__tests__/Login.test.jsx
import { render, screen, fireEvent } from 
  '@testing-library/react';
import userEvent from '@testing-library/user-event';
import { Provider } from 'react-redux';
import { BrowserRouter } from 'react-router-dom';
import store from '../../store';
import Login from '../Auth/Login';

describe('Login Component', () => {
  const renderLogin = () => {
    render(
      <Provider store={store}>
        <BrowserRouter>
          <Login />
        </BrowserRouter>
      </Provider>
    );
  };

  test('renders login form', () => {
    renderLogin();
    expect(screen.getByText('VisioLearn Login')). 
      toBeInTheDocument();
  });

  test('submits form with valid credentials', 
    async () => {
    renderLogin();
    const emailInput = screen.getByPlaceholderText('Email');
    const passwordInput = 
      screen.getByPlaceholderText('Password');
    const submitButton = 
      screen.getByRole('button', { name: /login/i });

    await userEvent.type(emailInput, 
      'admin@visiolearn.org');
    await userEvent.type(passwordInput, 
      'SecurePass123!');
    await userEvent.click(submitButton);

    // Assert login success
  });

  test('shows error on invalid credentials', 
    async () => {
    renderLogin();
    const emailInput = screen.getByPlaceholderText('Email');
    const passwordInput = 
      screen.getByPlaceholderText('Password');
    const submitButton = 
      screen.getByRole('button', { name: /login/i });

    await userEvent.type(emailInput, 'wrong@email.com');
    await userEvent.type(passwordInput, 'wrongpass');
    await userEvent.click(submitButton);

    // Assert error message
  });
});
\end{lstlisting}

% ============================================================================
% SECTION 11: PERFORMANCE & OPTIMIZATION
% ============================================================================
\section{Performance Optimization}

\subsection{Code Splitting}

\begin{lstlisting}[language=javascript]
// App.jsx - Route-based code splitting
import { lazy, Suspense } from 'react';
import LoadingSpinner from './components/Common/LoadingSpinner';

const Login = lazy(() => import('./pages/Login'));
const Dashboard = lazy(() => import('./pages/Dashboard'));
const VoiceSession = lazy(() => 
  import('./pages/VoiceSession'));

export default function App() {
  return (
    <Suspense fallback={<LoadingSpinner />}>
      <Routes>
        <Route path="/login" element={<Login />} />
        <Route path="/dashboard" element={<Dashboard />} />
        <Route path="/voice-session" 
          element={<VoiceSession />} />
      </Routes>
    </Suspense>
  );
}
\end{lstlisting}

\subsection{Memoization}

\begin{lstlisting}[language=javascript]
// Components - Use React.memo for expensive renders
import { memo } from 'react';

const InteractionLog = memo(({ interactions }) => {
  return (
    <div className="log">
      {interactions.map(i => (
        <div key={i.id}>{i.transcript}</div>
      ))}
    </div>
  );
});

export default InteractionLog;
\end{lstlisting}

\subsection{Lazy Loading Images}

\begin{lstlisting}[language=javascript]
// Use native lazy loading
<img 
  src="lesson-cover.jpg" 
  alt="Lesson" 
  loading="lazy" 
/>
\end{lstlisting}

% ============================================================================
% SECTION 12: DEPLOYMENT
% ============================================================================
\section{Deployment}

\subsection{Build Configuration}

\begin{lstlisting}[language=bash]
# .env.production
REACT_APP_API_BASE_URL=https://visiolearn-backend.onrender.com/api/v1
REACT_APP_ENV=production
\end{lstlisting}

\subsection{Build & Deploy}

\begin{lstlisting}[language=bash]
# Build for production
npm run build

# Output: dist/ folder
# Deploy to: Vercel, Netlify, GitHub Pages, etc.
\end{lstlisting}

\subsection{Recommended Platforms}

\begin{table}[H]
\centering
\small
\begin{tabular}{|l|l|}
\hline
\textbf{Platform} & \textbf{Features} \\
\hline
Vercel & Optimal for React, auto-deploy from GitHub \\
Netlify & Great DX, build optimization \\
GitHub Pages & Free, simple static hosting \\
\hline
\end{tabular}
\end{table}

% ============================================================================
% APPENDIX
% ============================================================================
\appendix

\section{Quick Reference}

\subsection{Common Commands}

\begin{table}[H]
\centering
\small
\begin{tabular}{|l|l|}
\hline
\textbf{Command} & \textbf{Purpose} \\
\hline
\texttt{npm install} & Install dependencies \\
\texttt{npm start} & Run dev server \\
\texttt{npm run build} & Build for production \\
\texttt{npm test} & Run tests \\
\texttt{npm run lint} & Run ESLint \\
\hline
\end{tabular}
\end{table}

\subsection{Useful Packages}

\begin{table}[H]
\centering
\small
\begin{tabular}{|l|l|}
\hline
\textbf{Package} & \textbf{Purpose} \\
\hline
\texttt{axios} & HTTP client \\
\texttt{react-router-dom} & Routing \\
\texttt{@reduxjs/toolkit} & State management \\
\texttt{react-query} & Server state management \\
\texttt{framer-motion} & Animations \\
\texttt{date-fns} & Date formatting \\
\hline
\end{tabular}
\end{table}

\section{Debugging Tips}

\begin{itemize}
    \item \textbf{Redux DevTools}: Monitor state changes in real-time
    \item \textbf{Network Tab}: Check API requests and responses
    \item \textbf{React DevTools}: Inspect component tree
    \item \textbf{Console}: Check for errors and warnings
    \item \textbf{Breakpoints}: Use debugger in VS Code
\end{itemize}

\section{Resources}

\begin{itemize}
    \item React Documentation: \url{https://react.dev}
    \item Redux Toolkit: \url{https://redux-toolkit.js.org}
    \item Axios: \url{https://axios-http.com}
    \item React Router: \url{https://reactrouter.com}
    \item Backend API: \url{https://visiolearn-backend.onrender.com/docs}
\end{itemize}

\end{document}
